\chapter{Lazy Evaluation, Graph Reduction, and Runtime Systems}
\label{ch:lazy-functional-runtime}

Delaying an expression is easy: retain its syntax and environment until a value is needed.
Sharing the result safely is harder.
The runtime must distinguish an unevaluated thunk, a thunk currently evaluating, a completed value, and a failure policy.
It must also avoid retaining dead environments after an update and distinguish productive cyclic data from circular demand.

\section{Represent suspended computation explicitly}
\label{sec:ch92-thunks}
A \term{thunk} stores a delayed expression and the lexical environment needed to evaluate it.
The call-by-need machine binds an argument parameter to a thunk instead of eagerly computing the argument.
Variable lookup forces that thunk; repeated lookup reuses a stored result after successful evaluation.
Let-bound values use the same mechanism, so two uses of one binding share work.

For let x=18+5 in x+x, the first x evaluates the delayed addition to 23 and updates its thunk.
The second x returns the cached 23.
The final answer is 46, and the lab's first-evaluation counter is one.
Table~\ref{tab:ch92-sharing} traces these events.
The counter counts newly evaluated thunks, not every source arithmetic operation or machine instruction.

\begin{table}[htbp]
\caption{Call-by-need updates one shared suspension after its first demand.}
\label{tab:ch92-sharing}
\centering
\begin{tabular}{p{0.20\linewidth}p{0.31\linewidth}p{0.34\linewidth}}
\toprule
Event & Thunk state & Returned value \\
\midrule
bind x & pending expression 18+5 & none yet \\
first use of x & busy, then done & 23 \\
second use of x & done & cached 23 \\
outer addition & done & 46 \\
\bottomrule
\end{tabular}
\end{table}

The thunk environment must come from the argument's call site, not from the function's definition site.
The function body itself still uses the closure's definition-site environment.
Confusing these two environments yields subtle scope errors that a first-order arithmetic test never reveals.
Nested closures and shadowing tests exercise both paths in the companion.

\section{Force with update and black-hole detection}
\label{sec:ch92-force}
Algorithm~\ref{alg:ch92-force} is the implemented forcing protocol.
Marking busy before evaluation detects reentry into the same unfinished computation.
The lab raises a black-hole error in that situation.
After successful evaluation it clears the retained expression and environment, reducing accidental retention of otherwise dead inputs.

\begin{algorithm}[htbp]
\caption{Force and update a single-threaded call-by-need thunk.}
\label{alg:ch92-force}
\begin{pseudocode}
\Require live thunk; serialized evaluation; explicit failure policy
\Ensure cached value on success or propagated failure without a false done state
\If{state is done} \State return cached value \EndIf
\If{state is busy} \State report circular demand \EndIf
\State mark busy and evaluate stored expression in stored environment
\If{evaluation fails} \State reset to pending and propagate failure \EndIf
\State store value; mark done; release expression/environment references
\State return cached value
\end{pseudocode}
\end{algorithm}

Successful repeated access costs constant thunk bookkeeping plus dispatch in this model.
The first force includes arbitrary expression evaluation and environment-copy costs.
The failure policy resets to pending rather than caching an exception; a retry can reevaluate the expression if a caller catches the host-side failure.
The source core has no exception handler, so the public CLI simply terminates with failure.
A language with catchable exceptions must specify whether failed thunks memoize the exception and its identity.

Figure~\ref{fig:ch92-states} shows the transitions.
In a multithreaded runtime, busy can mean another thread is computing the value rather than a true self-dependent black hole.
Ownership identity, wait queues, atomic publication, and cancellation then become necessary.
The single-threaded lab does not claim those concurrent semantics.

\begin{figure}[htbp]
\centering
\begin{tikzpicture}[node distance=19mm]
\node[basic block] (p) {pending};\node[basic block,right=of p] (b) {busy};
\node[basic block,right=of b] (d) {done};
\draw[cfg edge] (p)--(b);\draw[cfg edge] (b)--node[above,font=\small]{success}(d);
\draw[cfg edge,dashed] (b.south) to[bend left=35] node[below,font=\small]{failure resets}(p.south);
\end{tikzpicture}
\caption{Only successful evaluation publishes a done value in the reference runtime.}
\label{fig:ch92-states}
\end{figure}

\section{Observe weak-head structure before traversing data}
\label{sec:ch92-whnf}
\term{Weak-head normal form} exposes the outer constructor or lambda without necessarily normalizing its fields or body.
The lazy cons operation stores head and tail as thunks and returns the constructor pair immediately.
A list case forces the scrutinee only enough to distinguish Nil from Cons, then binds the existing field thunks.
Selecting an empty branch does not evaluate the cons branch, and selecting a head does not require traversing the tail.

This supports productive infinite structures in a richer recursive-data language, though the case-study surface restricts letrec to functions.
The distinction remains visible in finite examples with divergent field computations.
An eager conversion of a lazy list to a host list would force every tail and change both termination and memory behavior.
The CLI avoids that problem by printing scalar results only.

Figure~\ref{fig:ch92-list} shows one exposed list constructor.
The collector traces both field thunks and their environments while they remain reachable.
After a field is forced and updated, its old environment may become unreachable and reclaimable.
Whether that actually reduces memory depends on other live references, not solely on the update operation.

\begin{figure}[htbp]
\centering
\begin{tikzpicture}[node distance=14mm and 16mm]
\node[basic block] (c) {Cons: outer value available};
\node[basic block,below left=of c] (h) {head thunk};
\node[basic block,below right=of c] (t) {tail thunk};
\draw[cfg edge] (c)--(h);\draw[cfg edge] (c)--(t);
\end{tikzpicture}
\caption{Exposing a lazy constructor does not require evaluating either field.}
\label{fig:ch92-list}
\end{figure}

\section{Map graph reduction to runtime entry conventions}
\label{sec:ch92-entry}
An implementation can represent applications as graph nodes and reduce a spine until enough arguments reach a function entry.
Alternatively, compiled closures can enter native code directly with explicit argument and update frames.
The STG design separates function closures, constructor closures, and updateable suspended computations and gives a practical mapping to stock hardware~\cite{peytonjones-1992-stg}.
The lab's Python closure interpreter is a reference model, not an STG implementation.

Partial application needs a value retaining supplied arguments until enough arguments arrive.
Overapplication must handle a returned function receiving remaining arguments.
Arity information, entry points, and specialized calling paths reduce overhead, but every path must agree on which arguments are already evaluated and which are pointers to closures.
A backend also needs stack maps and update-frame metadata for GC and debugging.

Table~\ref{tab:ch92-runtime} distinguishes these runtime objects.
Treating a constructor as a zero-argument function can simplify one abstract presentation but can hide dispatch and allocation costs in a practical runtime.
Likewise, a thunk that returns a function still needs its own update protocol if it is shared.
The representation of the resulting function is not the representation of the suspended computation that produced it.

\begin{table}[htbp]
\caption{Lazy runtime objects have different entry and tracing obligations.}
\label{tab:ch92-runtime}
\centering
\begin{tabular}{p{0.20\linewidth}p{0.31\linewidth}p{0.34\linewidth}}
\toprule
Object & Entry behavior & Retained state \\
\midrule
function closure & consume argument(s), execute body & code and captures \\
thunk & evaluate once, update on success & expression/code and environment \\
constructor & expose tag and fields & field references \\
partial application & accumulate missing arguments & function and supplied values \\
update frame & install result after evaluation & target thunk and continuation \\
\bottomrule
\end{tabular}
\end{table}

\section{Analyze space, not only repeated work}
\label{sec:ch92-space}
Sharing can save computation and increase residency.
A thunk retaining an entire environment may keep a large input alive even if its eventual result is small.
A long chain of delayed accumulator updates can create a space leak where a strict accumulator would stay compact.
Conversely, forcing everything early destroys the ability to avoid unused divergent arguments.
Demand analysis and strictness information guide selective evaluation rather than replacing lazy semantics globally.

Algorithm~\ref{alg:ch92-update} states a runtime memory-accounting discipline for thunk transitions.
It is a design checklist, not a collector implemented by the Python lab.
Tracing cost depends on live closure graph size, including retained environments and update frames.
Cycles require a collector capable of reclaiming unreachable cycles or an explicitly different ownership scheme.

\begin{algorithm}[htbp]
\caption{Account for roots and obsolete references across a thunk update.}
\label{alg:ch92-update}
\begin{pseudocode}
\Require runtime GC protocol, active update target, result closure/value
\Ensure no live reference is lost and obsolete captures are not retained unnecessarily
\State keep update target and evaluation continuation visible as roots
\State publish the result through the runtime's required write barrier
\State change thunk state consistently with result visibility
\State remove obsolete expression and environment references
\State retire update-frame roots only after publication is complete
\end{pseudocode}
\end{algorithm}

For a concurrent or generational collector, publication and field clearing must follow the collector's barrier and memory-ordering rules.
The single-threaded Python object assignments rely on the host runtime's memory management and do not demonstrate a new barrier implementation.
Performance evaluation should report allocation and maximum live memory as well as elapsed time and number of thunk evaluations.
Counting reductions alone can hide the cost of graph manipulation and retained data.

\section{Follow a demand through a list case}
\label{sec:ch92-demand-example}
Consider a list whose head is integer four and whose tail computation loops forever, then a case returning the head in the Cons branch.
The lazy constructor allocates two field thunks and exposes Cons without entering the tail.
The case binds those existing thunks, evaluates the head variable, and returns four.
The tail remains pending and may become unreachable when the case finishes.
The strict machine instead evaluates the tail while constructing the list and never reaches the case result.
This difference follows from constructor-field strictness as well as argument strategy.

Demand analysis can identify fields or arguments that every terminating path will evaluate, allowing some thunk allocation to be avoided under a proven transformation.
But a syntactically present variable use is not enough to establish unconditional demand.
A use in an unselected branch may never execute, and an earlier divergent computation can prevent a later use.
The analysis must overapproximate the operational cases relevant to the particular strictness theorem rather than turning every occurrence into an eager evaluation point.

Sharing also changes how this example consumes memory.
If another closure retains the list, the unused tail thunk remains reachable even after the head-selecting case returns.
If no root retains the list, the constructor and pending tail can be reclaimed together by the host collector.
The thunk state machine tells us whether work has happened; reachability tells us whether the representation must remain allocated.
Neither fact alone predicts residency, which is why lazy-runtime measurements need both evaluation and heap information.

\section{Exercises}
\label{sec:ch92-exercises}
\begin{enumerate}
\item\label{ex:ch92-share} Trace let x=18+5 in x+x, giving the cached result and number of first-time thunk evaluations.
\item\label{ex:ch92-blackhole} A pending thunk's expression immediately forces that same thunk.
Give the state sequence and failure point.
\item\label{ex:ch92-fields} Explain why selecting the head of a lazy Cons need not evaluate its tail and why eagerly converting the entire value to a host list changes the contract.
\item\label{ex:ch92-retention} A completed thunk still retains its original environment containing a large input.
Which references can be cleared, and why does that not prove the input is globally dead?
\end{enumerate}

\section{Further reading}
\label{sec:ch92-reading}
Peyton Jones develops closure entry, updates, and compiled non-strict execution in the STG machine~\cite{peytonjones-1992-stg}.
The garbage-collection handbook provides the broader tracing and reclamation context~\cite{jones-2023-garbage}.
